% Replacement for Section V through the end of the uploaded source.
% Load section5_layout_preamble.tex at the end of the preamble first.
% This file is a document fragment, not a standalone document.

\FloatBarrier
\Needspace{8\baselineskip}
\section{Performance Evaluation}
\label{sec:performance_evaluation}

\subsection{Experimental Setup and Evaluation Protocol}
\label{subsec:experimental_setup}

We evaluate HAN--PDQN in a multi-user LEO satellite edge
computing network comprising 66 satellites distributed
across six orbital planes. The orbital altitude and
inclination are 550~km and $53^\circ$, respectively.
Ground users are distributed around
$(39.9^\circ\mathrm{N},116.4^\circ\mathrm{E})$, and the
population varies from 20 to 40 users.
User locations and orbital starting times are resampled
at episode initialization. Each episode contains
512 slots, with a slot duration of 1~s.
Tables~\ref{tab:simulation_parameters} and~\ref{tab:training_parameters}
summarize the system and training parameters, respectively.

% Place the two parameter tables together, outside the text columns.
\begin{strip}
\noindent
\begin{minipage}[t]{0.485\textwidth}
\vspace{0pt}
\LEOTableSetup
\captionof{table}{System Parameters}
    \label{tab:simulation_parameters}
    \centering
    \footnotesize
    \renewcommand{\arraystretch}{1.08}
    \setlength{\tabcolsep}{4pt}
    \begin{tabular}{p{0.60\linewidth}p{0.32\linewidth}}
        \hline
        Parameter & Value \\
        \hline
        Orbital planes / satellites per plane & $6/11$ \\
        Orbital altitude / inclination & $550$ km / $53^\circ$ \\
        Number of users, $N$ & $20,25,30,35,40$ \\
        User deployment radius & $5^\circ$ \\
        Slot duration, $\Delta t$ & $1$ s \\
        Episode length & $512$ slots \\
        Carrier frequency, $f_c$ & $20$ GHz \\
        Uplink bandwidth per satellite, $B_s$ & $10$ MHz \\
        User / satellite transmit power & $24/40$ dBm \\
        User / satellite antenna gain & $38.5/34$ dBi \\
        Noise temperature / noise figure & $354.81$ K / $2$ dB \\
        Local processing rate, $f_i$ & $1$ Gcycle/s \\
        Total MEC processing rate, $F_s$ & $10$ Gcycle/s \\
        MEC capacity / concurrency limit & $6/2$ tasks \\
        Task arrival probability & $0.35$ per user-slot \\
        Minimum elevation, $\theta_{\min}$ & $10^\circ$ \\
        Visibility thresholds, $v_{\min},v_{\mathrm{trig}}$
        & $60$ s \\
        Target SNR threshold, $\Gamma_{\min}$ & $-5$ dB \\
        Association-change overhead, $\delta^{\mathrm{ho}}$
        & $0.6$ s \\
        Output-to-input size ratio, $\zeta$ & $0.1$ \\
        Amplifier efficiency / circuit power & $0.38$ / $0.05$ W \\
        Local energy coefficient, $\kappa_i$ & $10^{-27}$ \\
        Minimum effective offloading fraction & $0.05$ \\
        \hline
    \end{tabular}
\end{minipage}\hfill%
\begin{minipage}[t]{0.485\textwidth}
\vspace{0pt}
\LEOTableSetup
\captionof{table}{Algorithm and Training Parameters}
    \label{tab:training_parameters}
    \centering
    \footnotesize
    \renewcommand{\arraystretch}{1.08}
    \setlength{\tabcolsep}{4pt}
    \begin{tabular}{p{0.60\linewidth}p{0.32\linewidth}}
        \hline
        Parameter & Value \\
        \hline
        Interaction budget per run & $150{,}000$ steps \\
        Training seeds & $42$--$46$ \\
        Evaluation episodes per seed & $5$ \\
        Discount factor, $\gamma$ & $0.99$ \\
        HAN layers / attention heads & $2/4$ \\
        HAN hidden / output dimension & $64/64$ \\
        PDQN Q-network hidden widths & $256,128$ \\
        PDQN parameter-branch hidden widths & $128,64$ \\
        PDQN learning rate & $10^{-3}$ \\
        HAN--PDQN / PDQN replay batch size & $512/128$ \\
        PDQN replay capacity & $50{,}000$ \\
        PDQN replay warm-up & $1{,}000$ steps \\
        PDQN target-copy interval & $500$ updates \\
        PDQN exploration probability & $1.0\rightarrow0.02$ \\
        MAPPO learning rate & $10^{-4}$ \\
        MAPPO rollout length / update epochs & $1024/4$ \\
        MAPPO minibatch size & $512$ \\
        MADDPG actor / critic learning rate
        & $10^{-4}/(3\times10^{-4})$ \\
        MADDPG replay batch size & $128$ \\
        MADDPG target soft-update coefficient & $0.005$ \\
        \hline
    \end{tabular}
\end{minipage}
\end{strip}

Each user generates a task independently with probability
0.35 per slot. Tasks are classified as light, medium,
and heavy with probabilities 0.3, 0.5, and 0.2.
For light tasks, the input size, computation demand,
and deadline are uniformly sampled from
$[1,3]$~Mbit, $[0.3,0.8]$~Gcycle, and $[1.5,3]$~s,
respectively. The corresponding ranges are
$[5,10]$~Mbit, $[1.2,2.5]$~Gcycle, and $[2.5,5.5]$~s
for medium tasks, and
$[12,25]$~Mbit, $[3,7]$~Gcycle, and $[4,8]$~s
for heavy tasks.
Each MEC provides a total processing rate of
10~Gcycle/s through two 5-GHz cores.
The utility weights are
$(\omega_{\mathrm D},\omega_{\mathrm E},
\omega_{\mathrm I},\omega_{\mathrm F},\omega_{\mathrm L})
=(0.60,0.40,0.30,0.20,0.05)$,
with $E_{\mathrm{ref}}=1$~J.

We compare HAN--PDQN with six benchmarks:
\begin{itemize}
    \item \textit{HAN+MAPPO}: A graph-assisted multi-agent
    proximal policy optimization controller for joint
    satellite selection and offloading.

    \item \textit{MAPPO}: A multi-agent proximal policy
    optimization controller using raw user, task, and
    candidate-satellite observations.

    \item \textit{MADDPG}: A multi-agent deep deterministic
    policy gradient controller using raw observations
    to generate satellite and offloading decisions.

    \item \textit{PDQN}: A parameterized Q-learning
    controller using raw observations without HAN
    representations.

    \item \textit{Min-Distance}: A distance-based policy
    that requests the nearest admissible satellite
    under the handover conditions and requests full
    offloading. It retains the current association
    when no shorter admissible alternative is available.

    \item \textit{Full-Local}: A policy that processes
    all tasks locally and issues no handover requests.
\end{itemize}
All methods operate under the same task, communication,
MEC service, and admission models.

For each user population, the learning methods are trained
with seeds 42--46 and an interaction budget of
$1.5\times10^5$ environment steps per run.
Table~\ref{tab:training_parameters} reports the principal
training settings.
HAN--PDQN and PDQN use replay batch sizes of 512 and 128,
respectively. Periodic model evaluation is performed
every 25,000 environment steps for the PDQN and MAPPO
variants, and every ten episodes for MADDPG.
Checkpoints are selected according to their mean
evaluation return.

Each selected policy is evaluated over five episodes
per training seed, using matched scenario seeds across
methods. Performance metrics are first averaged over
these episodes and then summarized across the five
training seeds. Shaded regions and error bars represent
95\% bootstrap confidence intervals at the training-seed
level.

The reported metrics follow
Section~\ref{subsec:task_computation_model}.
Successful-task delay is averaged over successful tasks,
whereas the deadline violation rate is normalized by
the number of generated tasks.
Energy per successful task divides the total accounted
terminal energy by the successful-task count.
We additionally report cumulative blocked user-time,
total terminal energy per episode, and handovers per
user-minute to characterize service continuity and
operating costs.

\Needspace{7\baselineskip}
\subsection{Training Convergence and Hyperparameter Sensitivity}
\label{subsec:training_performance}

Fig.~\ref{fig:algorithm_convergence} compares the training
returns of the five learning methods with 30 users.
The episode return accumulates the user-averaged slot
reward over an episode, and each curve averages the
corresponding returns across training seeds.
The figure includes complete episodes within the
interaction budget.

\par\addvspace{6pt}
\noindent\begin{minipage}{\columnwidth}
\LEOFigureSetup
\centering
    \includegraphics[
        trim=2.5mm 5mm 0mm 1mm,
        clip,
        width=\linewidth
    ]{algorithm_convergence_u30.pdf}
    \captionof{figure}{Training returns with 30 users and 512 steps
    per episode. Curves show five-seed means with 95\%
    bootstrap confidence intervals. No smoothing is applied;
    MAPPO variants retain their two-episode logging intervals.}
    \label{fig:algorithm_convergence}
\end{minipage}
\par\addvspace{8pt}

HAN+MAPPO and MAPPO improve rapidly during the initial
training stage and subsequently approach lower return
levels. HAN--PDQN and PDQN exhibit a more pronounced
increase during the intermediate stage, after which
HAN--PDQN maintains the highest mean return.
Averaged over the final 20 complete episodes,
HAN--PDQN achieves a return of 87.78, compared with
76.85 for PDQN. The corresponding returns of HAN+MAPPO
and MAPPO are approximately 59.30 and 57.86.
MADDPG exhibits lower late-stage returns and greater
variation across seeds.
The comparison therefore favors HAN--PDQN in the
quality of the learned policy under the allocated
interaction budget.

We next examine the sensitivity of HAN--PDQN to its
learning rate and replay batch size in
Fig.~\ref{fig:hyperparameter_convergence}.
For the learning-rate experiment, the batch size is
fixed at 512 and the rate is selected from
$\{10^{-4},3\times10^{-4},10^{-3},3\times10^{-3}\}$.
For the batch-size experiment, the learning rate is
fixed at $10^{-3}$ and the batch size is selected from
$\{128,256,512,1024\}$.
Each setting uses five training seeds.

\par\addvspace{6pt}
\noindent\begin{minipage}{\columnwidth}
\LEOFigureSetup
\centering
    \includegraphics[
        trim=5.5mm 0.8mm 0mm 5mm,
        clip,
        width=\linewidth
    ]{convergence_u30.pdf}
    \captionof{figure}{Hyperparameter sensitivity with 30 users:
    (a) learning rate with batch size 512;
    (b) replay batch size with learning rate $10^{-3}$.
    Curves show five-seed means and 95\% bootstrap confidence
    intervals after a three-episode trailing mean.}
    \label{fig:hyperparameter_convergence}
\end{minipage}
\par\addvspace{8pt}

In Fig.~\ref{fig:hyperparameter_convergence}(a), learning
rates between $10^{-4}$ and $10^{-3}$ attain similar
late-stage returns. Increasing the rate to
$3\times10^{-3}$ produces a substantially lower mean
return and a wider confidence interval.
This behavior indicates sensitivity to aggressive
updates when the value estimates and continuous
action parameters evolve together.

Fig.~\ref{fig:hyperparameter_convergence}(b) shows that
batch sizes 128, 512, and 1024 reach comparable return
levels, whereas batch size 256 yields a lower mean
and greater variation across seeds.
The dependence on batch size is therefore nonmonotonic
over the examined settings.
The default combination of learning rate $10^{-3}$
and batch size 512 attains a high return with a
consistent five-seed training profile and is used
for the subsequent HAN--PDQN comparisons.

\Needspace{7\baselineskip}
\subsection{Task Performance under User Scaling}
\label{subsec:user_scaling_performance}

Fig.~\ref{fig:task_performance} evaluates task service
as the user population increases from 20 to 40.
The per-user arrival process remains unchanged, so
a larger population increases the aggregate demand
for the shared uplink and MEC resources.
HAN--PDQN achieves the lowest mean successful-task
delay, deadline violation rate, and energy per successful
task at every evaluated population.

\begin{strip}
\noindent\begin{minipage}{\textwidth}
\LEOFigureSetup
\centering
    \includegraphics[
        trim=1.8mm 3.3mm 1.4mm 1.8mm,
        clip,
        width=\linewidth
    ]{figure1_line_metrics.pdf}
    \captionof{figure}{Task performance versus the number of users:
    (a) average successful-task delay;
    (b) deadline violation rate;
    (c) terminal energy per successful task.
    Curves show five-training-seed means, and shaded regions
    indicate 95\% bootstrap confidence intervals.}
    \label{fig:task_performance}
\end{minipage}
\end{strip}

In Fig.~\ref{fig:task_performance}(a), the mean
successful-task delay of HAN--PDQN increases from
1554.9~ms at 20 users to 1906.7~ms at 40 users.
The corresponding PDQN values are 1608.2 and 1939.9~ms.
At 30 users, HAN--PDQN achieves 1779.7~ms, compared
with 1820.7~ms for PDQN and 2053.8~ms for HAN+MAPPO.
The increase with population is consistent with
greater contention for upload bandwidth and MEC
service. Throughout this range, the proposed controller
maintains a lower delay despite the growing demand.

The reliability difference is more pronounced in
Fig.~\ref{fig:task_performance}(b).
At 30 users, the deadline violation rate is 6.89\%
for HAN--PDQN and 9.47\% for PDQN, a reduction of
2.58 percentage points.
HAN+MAPPO and MAPPO yield 18.38\% and 18.83\%,
respectively, under the same population.
At 40 users, HAN--PDQN maintains a violation rate
of 10.99\%, compared with 13.08\% for PDQN,
24.24\% for HAN+MAPPO, and 37.47\% for Min-Distance.
These results show that the delay advantage is
accompanied by improved deadline compliance.

Fig.~\ref{fig:task_performance}(c) further shows that
HAN--PDQN uses less terminal energy per successful
task across all evaluated populations.
Its value increases from 0.948~J at 20 users to
1.491~J at 40 users, while PDQN increases from
1.040 to 1.537~J.
At 30 users, HAN--PDQN requires 1.242~J per successful
task, compared with 1.318~J for PDQN and 2.376~J
for Min-Distance.
The combination of lower delay, fewer deadline violations,
and lower energy per successful task indicates an
improvement across the three task-service dimensions.

Full-Local exhibits nearly constant delay and energy
per successful task as the population changes, consistent
with users operating independent local processors.
Its deadline violation rate nevertheless remains
approximately 49\%, and its energy per successful
task is approximately 4.1~J.
The joint consideration of delay and task outcomes
therefore reveals the service limitations of relying
entirely on local computation.

For an aggregate comparison, we first take the arithmetic
mean of each metric over the five user populations
and then compute its relative reduction against PDQN.
Under this convention, HAN--PDQN reduces successful-task
delay, deadline violation rate, and energy per successful
task by 2.16\%, 21.26\%, and 5.72\%, respectively.
The largest relative improvement occurs in deadline
compliance, while the delay and energy improvements
remain consistent across the evaluated populations.

\Needspace{7\baselineskip}
\subsection{Service Continuity and Operating Costs}
\label{subsec:service_operating_costs}

Fig.~\ref{fig:service_operating_costs} complements the
task-level results by examining service interruption,
aggregate terminal energy, and handover activity.
These metrics reveal how different association and
offloading policies balance connection availability
against the conditions under which tasks are executed.

\begin{strip}
\noindent\begin{minipage}{\textwidth}
\LEOFigureSetup
\centering
    \includegraphics[
        trim=3mm 3.3mm 1.4mm 1.8mm,
        clip,
        width=\linewidth
    ]{figure2_bar_metrics.pdf}
    \captionof{figure}{Service continuity and operating costs:
    (a) blocked user-time per episode;
    (b) total terminal energy per episode;
    (c) handovers per user-minute.
    Bars show five-training-seed means with 95\% bootstrap
    confidence intervals. Full-Local is omitted from (a)
    and performs no handovers in (c).}
    \label{fig:service_operating_costs}
\end{minipage}
\end{strip}

In Fig.~\ref{fig:service_operating_costs}(a), HAN--PDQN
has lower mean blocked user-time than PDQN at every
population. At 30 users, the values are 284.08 and
298.12~user-s per episode, respectively.
The cumulative quantity increases with population,
partly because more users contribute to the total
observation time.
Min-Distance yields a slightly lower value of
276.32~user-s at 30 users, but its deadline violation
rate in Fig.~\ref{fig:task_performance}(b) is much higher.
Connection availability therefore needs to be considered
together with the processing conditions and workload
assigned to the selected satellite.

Fig.~\ref{fig:service_operating_costs}(b) shows that
HAN--PDQN also maintains the lowest mean total
terminal energy over the evaluated populations.
At 30 users, it consumes 6.18~kJ per episode,
compared with 6.37~kJ for PDQN,
8.41~kJ for Min-Distance, and 11.26~kJ for Full-Local.
Total energy measures the system's accumulated
terminal expenditure, whereas
Fig.~\ref{fig:task_performance}(c) measures the
expenditure per successful task.
The improvement in both quantities shows that the
energy-efficiency gain accompanies a reduction in
aggregate terminal consumption.

The handover results in
Fig.~\ref{fig:service_operating_costs}(c) show a
different ordering among the methods.
At 30 users, HAN--PDQN performs 0.606 handovers per
user-minute, compared with 0.659 for PDQN,
0.192 for HAN+MAPPO, and 0.198 for Min-Distance.
Thus, HAN--PDQN obtains its task-performance advantage
with fewer handovers than PDQN, while remaining more
active in satellite reassociation than HAN+MAPPO
and Min-Distance.
The lower switching activity of the latter policies
is accompanied by higher task delay, deadline violations,
and energy per successful task.
This comparison reflects the role of the long-term
objective in balancing immediate switching costs
against subsequent task-service conditions.

Using the same across-population averaging convention
as in Section~\ref{subsec:user_scaling_performance},
HAN--PDQN reduces mean blocked user-time,
total terminal energy, and handover frequency relative
to PDQN by 7.85\%, 3.19\%, and 14.82\%, respectively.
These operating-cost reductions accompany the
task-level improvements, linking the quality of
completed service to the association and workload
decisions that sustain it.

\FloatBarrier

% Insert the Conclusion here when ready; keep the bibliography below it.
\bibliographystyle{IEEEtran}
\bibliography{sample}
\end{document}